% Part: computability
% Chapter: recursive-functions
% Section: examples

\documentclass[../../../include/open-logic-section]{subfiles}

\begin{document}

\olfileid{cmp}{rec}{exa}
\olsection{आदिम पुनरावर्ती फलनांची उदाहरणे}


आदिम पुनरावर्ती फलनांची काही उदाहरणे आपल्याकडे आधीच आहेत: बेरीज आणि
गुणाकार ही फलने~$\Add$ आणि $\Mult$. एकरूपता फलन $\fn{id}(x) = x$ हे
आदिम पुनरावर्ती आहे, कारण ते फक्त~$\Proj{1}{0}$ आहे. स्थिर फलने
$\fn{const}_n(x) = n$ ही आदिम पुनरावर्ती आहेत, कारण $\Zero$ आणि $\Succ$
यांपासून क्रमिक संयोजनाने ती परिभाषित करता येतात. आदिम पुनरावर्ती
व्याख्यांत स्थिरांचा वापर करायचा असेल तेव्हा हे उपयोगी पडते. उदाहरणार्थ,
फलन $f(x) = 2 \cdot x$ परिभाषित करायचे असेल, तर $\fn{const}_n(x)$ आणि
गुणाकार यांपासून संयोजन करून ते $f(x) =
\Mult(\fn{const}_2(x), \Proj{1}{0}(x))$ असे मिळवता येते. यापुढे आपण
ही युक्ती वापरू.

\begin{prop}
  घातांक फलन $\fn{exp}(x, y) = x^y$ हे आदिम पुनरावर्ती आहे.
\end{prop}

\begin{proof}
  आपण $\fn{exp}$ ची आदिम पुनरावर्ती व्याख्या अशी देऊ शकतो:
  \begin{align*}
    \fn{exp}(x, 0) & = 1\\
    \fn{exp}(x, y+1) & = \Mult(x, \fn{exp}(x,y)).
    \intertext{काटेकोरपणे पाहता ही आदिम पुनरावर्ती फलनांपासून केलेली
      पुनरावर्ती व्याख्या नाही. मात्र औपचारिक रूपात आपल्याकडे पुढील
      व्याख्या आहे:}
    \fn{exp}(x, 0) & = f(x)\\
    \fn{exp}(x, y+1) & = g(x, y, \fn{exp}(x,y)).
    \intertext{येथे}
    f(x) & = \Succ(\Zero(x)) = 1\\
    g(x, y, z) & = \Mult(\Proj{3}{0}(x, y, z), \Proj{3}{2}(x, y, z)) = x \cdot z
  \end{align*}
  आणि म्हणून $f$ व $g$ ही आदिम पुनरावर्ती फलनांपासून संयोजनाने
  परिभाषित केलेली आहेत.
\end{proof}

\begin{prop}
  पुढीलप्रमाणे परिभाषित केलेले पूर्ववर्ती फलन $\fn{pred}(y)$,
  \[
  \fn{pred}(y) = \begin{cases}
    0 & \text{जर $y=0$ असेल}\\
    y-1 & \text{अन्यथा}
  \end{cases}
  \]
  हे आदिम पुनरावर्ती आहे.
\end{prop}

\begin{proof}
 पुढील गोष्ट लक्षात घ्या:
 \begin{align*}
   \fn{pred}(0) & = 0 \text{ आणि}\\
   \fn{pred}(y+1) & = y.
 \end{align*}
 ही जवळजवळ आदिम पुनरावर्ती व्याख्या आहे. काटेकोरपणे पाहता ती आदिम
 पुनरावर्तनाने केलेल्या व्याख्येच्या आकृतिबंधात बसत नाही, कारण त्या
 आकृतिबंधासाठी किमान एक अतिरिक्त आदान~$x$ आवश्यक असते. तसेच
 $\fn{pred}(y+1)$ च्या व्याख्येत $\fn{pred}(y)$ प्रत्यक्षात वापरले जात
 नाही, हेही काहीसे वेगळे आहे. पण आपण प्रथम पुढीलप्रमाणे
 $\fn{pred}'(x, y)$ परिभाषित करू शकतो:
 \begin{align*}
   \fn{pred}'(x, 0) & = \Zero(x) = 0,\\
   \fn{pred}'(x, y+1) & = \Proj{3}{1}(x, y, \fn{pred'}(x, y)) = y.
 \end{align*}
आणि नंतर त्यापासून संयोजनाने $\fn{pred}$ परिभाषित करू शकतो; उदाहरणार्थ,
$\fn{pred}(x) = \fn{pred}'(\Zero(x), \Proj{1}{0}(x))$ असे.
\end{proof}

\begin{prop}
  क्रमगुणित फलन $\fn{fac}(x) = \fact{x} = 1 \cdot 2 \cdot 3
  \cdot \dots \cdot x$ हे आदिम पुनरावर्ती आहे.
\end{prop}

\begin{proof}
  त्याची स्पष्ट आदिम पुनरावर्ती व्याख्या अशी आहे:
  \begin{align*}
    \fn{fac}(0) &= 1\\
    \fn{fac}(y+1) & = \fn{fac}(y) \cdot (y+1).
    \intertext{औपचारिक रूपात आपल्याला प्रथम दोन-स्थानी फलन $h$ परिभाषित करावे लागते:}
    h(x, 0) & = \fn{const}_1(x)\\
    h(x, y+1) & = g(x, y, h(x, y))
    \intertext{येथे $g(x, y, z) = \Mult(\Proj{3}{2}(x, y, z),
      \Succ(\Proj{3}{1}(x, y, z)))$; आणि नंतर पुढीलप्रमाणे ठेवावे:}
    \fn{fac}(y) & = h(\Proj{1}{0}(y), \Proj{1}{0}(y)) = h(y,y).
  \end{align*}
  यापुढे आपण थोडी अधिक औपचारिक मोकळीक घेऊ आणि संयोजन व आदिम
  पुनरावर्तन यांच्या साहाय्याने केलेल्या पूर्ण औपचारिक व्याख्या प्रत्येक
  वेळी देणार नाही.
\end{proof}

\begin{prop}
  पुढीलप्रमाणे परिभाषित केलेली छाटलेली वजाबाकी, $x \tsub y$,
  \[
  x \tsub y = \begin{cases}
    0 & \text{जर $x < y$ असेल}\\
    x-y & \text{अन्यथा}
  \end{cases}
  \]
  ही आदिम पुनरावर्ती आहे.
\end{prop}

\begin{proof}
  आपल्याकडे पुढील समीकरणे आहेत:
  \begin{align*}
    x \tsub 0 & = x\\
    x \tsub (y+1) & = \fn{pred}(x \tsub y) 
  \end{align*}
\end{proof}

\begin{prop}
 $x$ आणि $y$ यांतील अंतर, $\left|x-y\right|$, हे आदिम पुनरावर्ती आहे.
\end{prop}

\begin{proof}
  आपल्याकडे $\left| x-y \right| = (x \tsub y) + (y \tsub x)$ आहे;
  म्हणून अंतर हे $+$ आणि $\tsub$ यांपासून संयोजनाने परिभाषित करता येते,
  आणि ही दोन्ही आदिम पुनरावर्ती आहेत.
\end{proof}

\begin{prop}
  $x$ आणि $y$ यांतील महत्तम, $\fn{max}(x,y)$, हे आदिम पुनरावर्ती आहे.
\end{prop}

\begin{proof}
  $+$ आणि $\tsub$ यांपासून संयोजनाने $\fn{max}(x,y)$ ची व्याख्या अशी
  देता येते:
  \[
  \fn{max}(x,y) \defis x + (y \tsub x).
  \]
  $x$ हे महत्तम असेल, म्हणजे $x \ge y$ असेल, तर $y \tsub x = 0$;
  म्हणून $x + (y \tsub x) = x + 0 = x$. $y$ हे महत्तम असेल, तर
  $y \tsub x = y - x$, आणि म्हणून $x + (y \tsub x) = x + (y - x) = y$.
\end{proof}

\begin{prop}
  \ollabel{prop:min-pr}
  $x$ आणि $y$ यांतील लघुतम, $\fn{min}(x,y)$, हे आदिम पुनरावर्ती आहे.
\end{prop}

\begin{proof}
  सराव.
\end{proof}

\begin{prob}
  \olref[cmp][rec][exa]{prop:min-pr} सिद्ध करा.
\end{prob}

\begin{prob}
पुढील फलन आदिम पुनरावर्ती आहे हे दाखवा:
\[
f(x, y) =
2^{(2^{\iddots^{2^{x}}})}\raisebox{1ex}{\bigg\rbrace}
\raisebox{1ex}{\text {$y$ $2$'s}}\]
\end{prob}

\begin{prob}
पूर्णांक भागाकार $d(x, y) = \lfloor x/y \rfloor$ (म्हणजे दशांश
बिंदूनंतर येणारे सर्व काही वगळलेला भागाकार) हा आदिम पुनरावर्ती आहे, हे
दाखवा. $y = 0$ असताना $d(x, y) = 0$ असे आपण ठरवतो. आदिम पुनरावर्तन
आणि संयोजन वापरून~$d$ ची स्पष्ट व्याख्या द्या.
\end{prob}

\begin{prop}
आदिम पुनरावर्ती फलनांचा संच पुढील दोन क्रियांखाली बंद आहे:
\begin{enumerate}
\item सांत बेरजा: $f(\vec x, z)$ हे आदिम पुनरावर्ती असेल, तर पुढील
फलनही आदिम पुनरावर्ती असते:
\[
g(\vec x, y) \defis \sum_{z = 0}^y f(\vec x, z).
\]
\item सांत गुणाकार: $f(\vec x, z)$ हे आदिम पुनरावर्ती असेल, तर पुढील
फलनही आदिम पुनरावर्ती असते:
\[
h(\vec x, y) \defis \prod_{z = 0}^y f(\vec x, z).
\]
\end{enumerate}
\end{prop}

\begin{proof}
उदाहरणार्थ, सांत बेरजा पुढील समीकरणांनी पुनरावर्ती रीतीने परिभाषित होतात:
\begin{align*}
  g(\vec x, 0) & = f(\vec x, 0)\\
  g(\vec x, y+1) & = g(\vec x, y) + f(\vec x, y+1).
\end{align*}
\end{proof}


\end{document}
